\documentclass[11pt,a4paper]{article}
\usepackage[T1]{fontenc}
\usepackage[utf8]{inputenc}
\usepackage[spanish,es-nodecimaldot]{babel}
\usepackage{lmodern}
\usepackage[margin=24mm]{geometry}
\usepackage{amsmath,amssymb,booktabs,array}
\usepackage{xcolor,tikz,fancyhdr}
\usetikzlibrary{arrows.meta,positioning,calc,babel}
\definecolor{ink}{HTML}{17324D}
\definecolor{accent}{HTML}{087E8B}
\definecolor{pale}{HTML}{EDF5F7}
\usepackage[colorlinks=true,linkcolor=accent,urlcolor=accent]{hyperref}
\hypersetup{pdftitle={EMLLM: redes en un lenguaje One/Magma},pdfauthor={Matías / SugarFreeDoNoSo}}
\pagestyle{fancy}
\fancyhf{}
\fancyhead[L]{\small\textcolor{ink}{EMLLM / Informe conceptual}}
\fancyhead[R]{\small One/Magma}
\fancyfoot[L]{\small Diseño propuesto, no resultados de ejecución}
\fancyfoot[R]{\thepage}
\setlength{\headheight}{14pt}
\setlength{\parindent}{0pt}
\setlength{\parskip}{7pt}
\setlength{\emergencystretch}{2em}
\tikzset{box/.style={draw=accent,rounded corners=3pt,fill=pale,align=center,minimum height=12mm,text width=32mm,font=\small},arr/.style={-{Latex[length=2mm]},thick,draw=ink},neuron/.style={circle,draw=accent,fill=pale,minimum size=10mm}}
\newcommand{\M}{\operatorname{M}}
\newcommand{\One}{\mathsf{One}}
\newcommand{\Magma}{\mathsf{Magma}}

\begin{document}
\thispagestyle{empty}
{\color{accent}\large NOTA TÉCNICA / 01}\par
\vspace{10mm}
{\color{ink}\Huge\bfseries EMLLM}\par
\vspace{4mm}
{\color{ink}\LARGE Redes neuronales expresadas\par únicamente con One y Magma}\par
\vspace{6mm}
{\large Matías / \texttt{SugarFreeDoNoSo}}\par
\textbf{Versión documental:} 0.1.0\par
\textbf{Estado:} propuesta de arquitectura e implementación inicial.

\vspace{5mm}
\section*{Resumen}
El proyecto propone representar una inferencia neuronal como un programa
simbólico cerrado con una constante, \(\One\), y una operación binaria,
\(\Magma\). La semántica adoptada es
\[
  \operatorname{eval}(\One)=1,\qquad
  \M(x,y)=e^x-\ln y,\quad y>0.
\]
Entradas, pesos, sesgos y operaciones de la red deben convertirse a ese
lenguaje antes de evaluar sus salidas. El primer objetivo es construir y
verificar esta representación en Rust, empezando por una red pequeña.

\begin{center}
\begin{tikzpicture}[node distance=8mm]
\node[box] (data) {Entradas y pesos\\racionales};
\node[box,right=of data] (expr) {Expansión\\One/Magma};
\node[box,right=of expr] (dag) {DAG completo\\sin evaluar};
\node[box,below=10mm of dag] (eval) {Evaluación\\numérica};
\draw[arr] (data)--(expr);
\draw[arr] (expr)--(dag);
\draw[arr] (dag)--(eval);
\end{tikzpicture}
\end{center}

\textbf{Decisión central.} No habrá hojas \texttt{f64} de parámetros ni
instrucciones primitivas de suma o producto en el programa ejecutable.
Los números de máquina pertenecen al evaluador y a la interfaz de construcción,
no al alfabeto del programa EML.

\textbf{Lectura de la evidencia.} Este informe describe una propuesta.
El repositorio todavía no contiene el motor Rust. El prototipo citado en
la conversación no está incorporado ni verificado aquí; sus métricas no
se presentan como resultados de EMLLM.

\vfill
\href{https://github.com/SugarFreeDoNoSo/emllm}{github.com/SugarFreeDoNoSo/emllm}
\newpage

\section{Lenguaje y semántica}
La gramática del programa es deliberadamente mínima:
\[
 T ::= \One\;\mid\;\Magma(T,T).
\]
La evaluación es parcial: un nodo no puede evaluarse en los reales si el
valor de su segundo operando no es estrictamente positivo. Además, una
implementación finita puede fallar por desbordamiento o pérdida de precisión.
La denominación EML se usa aquí para la fórmula anterior, en base natural;
no debe intercambiarse con una operación de base dos.

\subsection{Construcciones elementales}
Sin introducir nuevas constantes primitivas se obtienen, en aritmética real:
\begin{align*}
 t_1 &= \M(1,1)=e,\\
 t_2 &= \M(t_1,1)=e^e,\\
 t_3 &= \M(1,t_2)=e-\ln(e^e)=0.
\end{align*}
Los nombres \(t_i\) son referencias a expresiones, no constantes adicionales.
El mismo nodo \(\One\) puede compartirse entre todos los usos.

\begin{center}
\begin{tikzpicture}[node distance=10mm]
\node[box] (one) {\(\One\)\\valor: 1};
\node[box,right=of one] (e) {\(t_1=\M(1,1)\)\\valor: \(e\)};
\node[box,below=of e] (ee) {\(t_2=\M(t_1,1)\)\\valor: \(e^e\)};
\node[box,left=of ee] (zero) {\(t_3=\M(1,t_2)\)\\valor: 0};
\draw[arr] (one)--(e);
\draw[arr] (e)--(ee);
\draw[arr] (one)--(ee);
\draw[arr] (one)--(zero);
\draw[arr] (ee)--(zero);
\end{tikzpicture}
\end{center}
\textit{Las flechas indican dependencias, no una ejecución durante la construcción.}

\subsection{Macros y frontera numérica}
Suma, negación y producto serán constructores que \emph{expanden} a One/Magma.
Sus nombres pueden existir en la API, pero no como instrucciones ejecutables.
Cada identidad utilizada debe documentar su derivación y dominio.

La completitud se toma como premisa de trabajo, no como un teorema demostrado
en este informe. Tampoco implica que todo real admita una representación
finita exacta, ni determina un algoritmo eficiente de síntesis. La primera
interfaz se limita a racionales acotados. Como mínimo admitirá \(p/q\), con
\(-2\leq p\leq2\) y \(1\leq q\leq2\).
\newpage

\section{Arquitectura propuesta}
\subsection{Construir primero; evaluar después}
Una arena almacena nodos inmutables y referencias a sus hijos. La identidad
es estructural: dos nodos Magma se comparten únicamente si tienen los mismos
hijos en el mismo orden. Coincidir numéricamente no los vuelve idénticos.

\begin{center}
\begin{tikzpicture}[node distance=8mm]
\node[box,text width=43mm] (front) {Interfaz de construcción\\racionales y dimensiones};
\node[box,text width=43mm,right=12mm of front] (macro) {Síntesis y macros\\solo producen One/Magma};
\node[box,text width=43mm,below=of macro] (graph) {DAG inmutable\\raíces y referencias};
\node[box,text width=43mm,left=12mm of graph] (valid) {Validación y planificación\\dominio estructural y límites};
\node[box,text width=43mm,below=of valid] (runtime) {Evaluador topológico\\\texttt{f64}, exp y ln};
\node[box,text width=43mm,right=12mm of runtime] (output) {Salida o error explícito\\nodo y categoría};
\draw[arr] (front)--(macro);
\draw[arr] (macro)--(graph);
\draw[arr] (graph)--(valid);
\draw[arr] (valid)--(runtime);
\draw[arr] (runtime)--(output);
\end{tikzpicture}
\end{center}

Construir el grafo completo no elimina sus dependencias numéricas. Durante
la evaluación, cada nodo sigue necesitando los valores de sus hijos. La
separación permite inspección, compartición y futuras optimizaciones sin
confundirlas con una aceleración ya demostrada.

\subsection{Entradas y pesos también son expresiones}
La red resultante es cerrada: no recibe valores mediante nodos Input o
Parameter. Para otra entrada se construyen nuevas raíces, compartiendo
subexpresiones inmutables cuando corresponda. Este coste es una consecuencia
explícita de la elección One/Magma, no algo que deba ocultarse mediante enlaces.

\subsection{Fallos observables}
El constructor debe rechazar dimensiones y referencias inválidas, ciclos y
presupuestos excedidos. El evaluador debe rechazar argumentos logarítmicos
no positivos, resultados no finitos y desbordamientos. No se usarán ajustes
silenciosos del dominio ni sustituciones por resultados convencionales.

La primera versión conservará los valores intermedios para facilitar
diagnósticos. La reutilización de temporales y la optimización algebraica
general quedan como posibles evoluciones, no como requisitos iniciales.
\newpage

\section{Demostración neuronal y aceptación}
Se propone una red densa 2--2--1, suficientemente pequeña para inspeccionar
su expansión completa:
\[
 h=g(W_1x+b_1),\qquad y=W_oh+b_o,\qquad g(z)=e^{-z^2}.
\]
Esta notación describe la función; la ejecución debe utilizar exclusivamente
su expansión One/Magma.

\begin{center}
\begin{tikzpicture}
\node[neuron] (x1) at (0,1) {\(x_1\)};
\node[neuron] (x2) at (0,-1) {\(x_2\)};
\node[neuron] (h1) at (4,1) {\(h_1\)};
\node[neuron] (h2) at (4,-1) {\(h_2\)};
\node[neuron] (y) at (8,0) {\(y\)};
\foreach \i in {1,2} {\foreach \j in {1,2} {\draw[arr] (x\i)--(h\j);}}
\draw[arr] (h1)--(y); \draw[arr] (h2)--(y);
\node[font=\small] at (0,-1.9) {Entradas};
\node[font=\small] at (4,-1.9) {Activación gaussiana};
\node[font=\small] at (8,-1.9) {Salida afín};
\end{tikzpicture}
\end{center}

Las fixtures previstas fijan:
\[
 W_1=\begin{pmatrix}\frac12&-\frac12\\0&\frac12\end{pmatrix},\quad
 b_1=\begin{pmatrix}0\\-\frac12\end{pmatrix},\quad
 W_o=\begin{pmatrix}\frac12&-\frac12\end{pmatrix},\quad b_o=0.
\]
Se probarán las entradas \((0,0)\), \((\frac12,-\frac12)\) y
\((-\frac12,\frac12)\). Una referencia aritmética independiente calculará
la misma función; no formará parte del evaluador EML.

\begin{center}
\begin{tabular}{@{}p{88mm}r@{}}
\toprule
\textbf{Comprobación prevista} & \textbf{Error absoluto máximo}\\
\midrule
\(\M(1,1)\) frente a \(e\) & \(10^{-12}\)\\
Racionales del conjunto de aceptación & \(10^{-9}\)\\
Macros en \(\{-\frac12,0,\frac12\}\) & \(10^{-8}\)\\
Salida de las tres fixtures neuronales & \(10^{-6}\)\\
\bottomrule
\end{tabular}
\end{center}

\textbf{Estos son objetivos de aceptación, no medidas obtenidas.}
También se comprobará que todos los nodos alcanzables son One/Magma y se
informarán tamaño del grafo, presupuestos y error observado. Una discrepancia
numérica debe bloquear la aceptación, no provocar una relajación silenciosa
de las pruebas.
\newpage

\section{Alcance, evidencia y publicación}
\subsection{Qué aporta esta primera etapa}
El entregable de implementación previsto es un núcleo Rust con construcción,
evaluación, síntesis acotada y demostración neuronal. Entrenamiento,
destilación, Transformer, importación GGUF y síntesis de pesos arbitrarios
quedan fuera de esta primera etapa. Su incorporación requerirá propuestas
posteriores y evidencia propia.

El prototipo descrito en la conversación utilizaba un MLP y un profesor
Markoviano experimental. Conservaba Add, Mul, Neg y parámetros numéricos;
por tanto, no constituye una demostración de una red puramente One/Magma.
Los resultados de Python/PyTorch allí comunicados tampoco verifican una
implementación Rust que aún no está presente en este repositorio.

\subsection{Riesgos técnicos que guían la implementación}
\begin{itemize}
\item \textbf{Estabilidad:} exponenciales, logaritmos y cancelación pueden
      amplificar errores. Por ejemplo, la identidad que construye cero no
      garantiza un cero exacto en coma flotante.
\item \textbf{Tamaño:} una expansión puede crecer mucho aunque la red sea
      pequeña. Se usarán límites explícitos y compartición estructural.
\item \textbf{Coste:} representar una multiplicación con EML no implica
      que sea más barata que una instrucción nativa.
\item \textbf{Escalado:} superar las fixtures no demuestra viabilidad para
      un LLM. Primero se necesita una ejecución pequeña y reproducible.
\end{itemize}

\subsection{Documento versionado y release automatizada}
La fuente de este informe vive en \texttt{docs/report/emllm.tex}. GitHub
Actions compila el documento con LaTeX y conserva el PDF como artefacto.
Los tags \texttt{report-v*} publican además una release con el PDF y su
checksum SHA-256. El tag vincula cada publicación a una revisión concreta.

\begin{center}
\begin{tikzpicture}[node distance=7mm]
\node[box,text width=29mm] (git) {Commit y tag\\\texttt{report-v*}};
\node[box,text width=29mm,right=of git] (ci) {GitHub Actions\\LaTeX};
\node[box,text width=29mm,right=of ci] (pdf) {PDF verificado\\SHA-256};
\node[box,text width=29mm,right=of pdf] (rel) {GitHub\\Release};
\draw[arr] (git)--(ci); \draw[arr] (ci)--(pdf); \draw[arr] (pdf)--(rel);
\end{tikzpicture}
\end{center}

\subsection{Fuentes y trazabilidad}
La fuente normativa es la propuesta OpenSpec
\texttt{implement-one-magma-eml}: propuesta, diseño, especificaciones de
expresiones e inferencia, y tareas. Se encuentra en
\href{https://github.com/SugarFreeDoNoSo/emllm/tree/main/openspec/changes/implement-one-magma-eml}{el repositorio del proyecto}.
Este informe resume esas decisiones; no sustituye sus criterios de aceptación
ni presenta una demostración formal de completitud.
\end{document}
